\documentclass[11pt,a4paper]{article}
\usepackage[T1]{fontenc}
\usepackage{isabelle,isabellesym}
\usepackage{amsmath,amssymb}
% Keep short formal statements together without another package dependency.
\newcommand{\ProofNeedspace}[1]{\par\begingroup
  \dimen0=#1\relax\dimen2=\pagegoal\advance\dimen2 by -\pagetotal
  \ifdim\dimen0>\dimen2\newpage\fi\endgroup}
\usepackage{pdfsetup}
\urlstyle{rm}
\isabellestyle{it}
\begin{document}
\title{Monotone--Light Factorization of Compact Hausdorff Maps}
\author{Arthur Freitas Ramos \\ David Barros Hulak \\
Ruy Jose Guerra Barretto de Queiroz}
\maketitle
\begin{abstract}
Every continuous map from a compact Hausdorff space to a Hausdorff space
factors through a compact Hausdorff space as a surjective monotone map
followed by a light map. The middle space is the quotient by connected
components of point fibres. We prove the factorization and its uniqueness
up to a homeomorphism commuting with the factor maps.
\end{abstract}
\tableofcontents
\section{Introduction}
For a continuous map $f:X\to Y$, a monotone map has connected point fibres,
and a light map has totally disconnected point fibres. These notions isolate
two different behaviours: the collapse of connected subsets, and maps whose
fibres have no nontrivial connected subset. The classical monotone--light
factorization separates these behaviours. Eilenberg's original result treats
compact metric spaces~\cite{Eilenberg1934}; the compact Hausdorff formulation
and its categorical interpretation are discussed by
Carboni, Janelidze, Kelly and Par\'e~\cite{Carboni1997}, Sections 7.3--7.5.
Michael~\cite{Michael1964}, Section 2, pages 5--6, gives a direct account
of fibre components, lightness and Hausdorff separation.

Our statement assumes that $X$ is compact Hausdorff and $Y$ is Hausdorff.
The map $f$ need not be surjective, and $Y$ need not be compact. The empty
space is allowed. The result constructs a middle space $M$, a continuous
surjective quotient $q:X\to M$, and a continuous map $l:M\to Y$, such that
$f=l\circ q$ on the carrier of $X$. The fibres of $q$ are precisely the
connected components of the fibres of $f$. Consequently $q$ is monotone
and $l$ is light.

This is a formalization of a classical theorem, with no claim of new
mathematical content. A targeted search of HOL-Analysis and relevant AFP
topology entries found no existing formalization of this factorization.
The general quotient topology is imported from the AFP entry
\emph{Topological Groups}~\cite{TopologicalGroups}; the development uses
HOL-Analysis's existing connected components, quasicomponents, monotone
maps, separation theorems and quotient lifting theorem.

Codex assisted with literature searches, proof planning, Isabelle proof
engineering, documentation and validation. An independent AI mathematical
review checked the argument and the generality of its statements. This
disclosure does not assign individual duties to the established human
author team or claim human review or endorsement.

\section{Proof and formal representation}
The carrier of the middle space is the set of fibre components. The
projection sends $x$ to the component of $x$ in $f^{-1}\{f(x)\}$, with the
quotient topology. Equal projection values characterize membership in the
same fibre component. Isabelle represents spaces as topology values with
an explicit carrier, so factor identities and uniqueness are asserted on
carriers; total functions may have irrelevant values outside them.

The central point is that $q$ is closed. Given a closed $A\subseteq X$ and
a point $x$ whose fibre component avoids $A$, work inside the compact
Hausdorff fibre $F=f^{-1}\{f(x)\}$. Components and quasicomponents coincide
there. A clopen subset $P$ of $F$ contains the component of $x$ and avoids
$A\cap F$. Normality of $X$ gives disjoint open neighbourhoods $U,V$ of
$P$ and $A\cup(F\setminus P)$ respectively. Since $f$ is closed, an open
neighbourhood $N$ of $f(x)$ has $f^{-1}N\subseteq U\cup V$. A connected
fibre component through $U\cap f^{-1}N$ cannot meet $V$, so this open set
is saturated and avoids $A$. The saturation of every closed $A$ is
therefore closed. The quotient criterion proves that $q$ is closed.

Compactness of $M$ follows from the continuous surjection $q$.
A closed continuous image of a normal Hausdorff space is Hausdorff.
The quotient lifting theorem now induces $l$. For a connected subset $C$
of an $l$-fibre, the closed monotone map $q$ has connected preimage
$q^{-1}C$. That preimage is contained in an $f$-fibre and hence in one
component. Surjectivity of $q$ forces $C$ to have at most one point,
which proves lightness.

For any alternative continuous monotone--light factorization through a
Hausdorff middle space, with a surjective first map, its first-map fibres equal the fibre components
of $f$: monotonicity gives one inclusion, and lightness of the image
of a connected component gives the other. Quotient lifting produces
the comparison map. It is a continuous bijection from compact $M$
to a Hausdorff space, hence a homeomorphism. Surjectivity of $q$
makes that comparison unique on the carrier.

All new definitions, statements and complete proofs follow in the formal
development.
\input{session}
\bibliographystyle{abbrv}
\bibliography{root}
\end{document}
